\subsection{Fixed-design least squares with a quantized right-hand side}
\label{sec:revision-ls}
The covariance formula also predicts errors in coefficients of a prescribed
linear approximation. Let $A\in\mathbb R^{n\times p}$ be deterministic and
have full column rank, and let $b_i$ be i.i.d. observations. Define
$\theta=A^\dagger b$ and $\widetilde\theta=A^\dagger Q(b)$, where the
right-hand side is quantized as one block and
$A^\dagger=(A^TA)^{-1}A^T$. These are projections of the same observations;
we measure the additional coefficient error due to data compression.

\begin{corollary}[Coefficient error and the constant direction]
\label{cor:revision-ls}
Put $g=A^\dagger\mathbf1$ and $\delta\theta=\widetilde\theta-\theta$. Then
\begin{align}
 \mathbb E[\delta\theta]&=\mu g,\qquad
 \operatorname{Cov}(\delta\theta)
 =(v-\kappa)(A^TA)^{-1}+\kappa gg^T,\label{eq:revision-ls-cov}\\
 \mathbb E[\|\delta\theta\|_2^2]
 &=(v-\kappa)\operatorname{tr}((A^TA)^{-1})
       +(\kappa+\mu^2)\|g\|_2^2.\label{eq:revision-ls-mse}
\end{align}
If $A=[\mathbf1,Z]$ has full column rank, then $g=(1,0,\ldots,0)^T$.
If instead $A^T\mathbf1=0$, then $g=0$.
\end{corollary}
\begin{proof}
Since $\delta\theta=A^\dagger e$, multiply
\eqref{eq:revision-spectrum} on both sides by $A^\dagger$ and its transpose,
and use $A^\dagger(A^\dagger)^T=(A^TA)^{-1}$. The second identity is the
trace of the covariance plus the squared mean. For the stated designs,
$A(1,0,\ldots,0)^T=\mathbf1$ or $A^T\mathbf1=0$, respectively.
\end{proof}

Thus the inverse Gram matrix controls ordinary geometric amplification,
while $g$ determines how the shared component and bias enter the fitted
coefficients. An intercept confines this component to the intercept
coefficient; omitting the intercept can transmit it to a slope. The centered
slope still has variance proportional to $v-\kappa$, rather than $v$.

We use $n=128$ and $t_i=(i-1)/(n-1)$. The first design fits
$\alpha+\beta(t_i-1/2)$, and the second fits $\gamma t_i$ through the
origin. Inputs are either $\operatorname{Exp}(1)$ or
$\exp(N(0,1))$; there is no deterministic trend added to these i.i.d.
right-hand sides. Each law uses $80\,000$ independent blocks and $400$
block bootstrap replicates. At $L=2$, the intercept biases are $-0.04222$
and $-0.23654$, respectively. For lognormal inputs, retaining the exact
marginal $\mu,v$ but setting $\kappa=0$ underpredicts the intercept and
through-origin slope variances by factors $9.88$ and $7.62$,
while overpredicting the centered slope variance by $7.5\%$.
The mean, variance, and mean-square-error intervals cover $6$,
$5$, and $5$ of $6$ predictions, respectively.
The retained variance-interval misses are: Exponential centered slope (exact ratio $0.992981$, interval $[0.993090,1.016730]$). Full coefficient results are in the package.

For an accuracy decision, $\alpha=\overline b$ has predicted mean-square
compression error
$[v+(n-1)\kappa]/n+\mu^2$. Table~\ref{tab:revision-ls} uses the same
input blocks at $L=2,3,4,5$. Besides the bias-aware independent baseline
$v/n+\mu^2$, it reports the classical independent, zero-mean uniform-cell
prediction $\mathbb E[s^2]/(12n)$. We choose an illustrative RMS tolerance
$0.10$ after inspecting the initial experiment; this is a transparent
design example, not a preregistered threshold or a precision optimizer.

\begin{table}[tb]
\centering
\caption{RMS compression error in the fitted intercept $\alpha$.
``Ind.'' retains the exact bias and marginal variance but drops pair
covariance; ``Unif.'' assumes independent, zero-mean uniform-cell errors.
The last column gives $95\%$ block bootstrap intervals for the measured
RMS. All tested precisions are retained.}
\label{tab:revision-ls}
\begingroup\footnotesize\setlength{\tabcolsep}{3pt}
\begin{tabular}{llrrrr}
\toprule
Law & $L$ & Exact & Ind. & Unif. & MC RMS interval\\
\midrule
Exp(1) & 2 & 0.0548 & 0.0493 & 0.0262 & [0.0545, 0.0552] \\
Exp(1) & 3 & 0.0182 & 0.0170 & 0.0131 & [0.0181, 0.0183] \\
Exp(1) & 4 & 0.0073 & 0.0071 & 0.0066 & [0.0072, 0.0073] \\
Exp(1) & 5 & 0.0034 & 0.0033 & 0.0033 & [0.0034, 0.0034] \\
Lognormal & 2 & 0.3138 & 0.2455 & 0.0782 & [0.3124, 0.3165] \\
Lognormal & 3 & 0.1211 & 0.0760 & 0.0391 & [0.1204, 0.1236] \\
Lognormal & 4 & 0.0395 & 0.0229 & 0.0195 & [0.0390, 0.0413] \\
Lognormal & 5 & 0.0129 & 0.0100 & 0.0098 & [0.0123, 0.0132] \\
\bottomrule
\end{tabular}\endgroup
\end{table}

For lognormal inputs, among the tested precisions, the exact prediction
selects $L=4$, whereas the bias-aware independent model selects
$L=3$ and the uniform-cell model selects $L=2$.
Monte Carlo shows that the exact selection meets the tolerance (measured RMS $0.04015$); the bias-aware independent selection fails the tolerance (measured RMS $0.12199$); the uniform-cell selection fails the tolerance (measured RMS $0.31446$). For exponential inputs all three models select
$L=2$; the measured RMS at this precision is $0.05481$.
At fixed block length, increasing
$L$ by one increases a sufficient signed integer coefficient width by
one bit while preserving one shared exponent; this is a model-level
storage tradeoff, not a hardware format comparison. Only the right-hand
side is quantized. The designs are fixed and their QR pseudoinverses are
applied to the residual in binary64, isolating input compression from
subsequent solver arithmetic and from non-i.i.d. regression models.
